\documentclass{article}
\usepackage[dvips]{graphicx}
\usepackage{a4wide}
\usepackage{amsmath}
\usepackage{euscript}
\usepackage{amssymb}
\usepackage{amsthm}
\usepackage{amsopn}

\theoremstyle{definition}
\newtheorem*{definition}{Definition}
\newtheorem{theorem}{Theorem}

\renewcommand{\qed}{\blacksquare}
\newcommand{\powset}{{\cal P}}

\newcommand{\vv}{\ensuremath{\vec{v}}}
\newcommand{\vu}{\ensuremath{\vec{u}}}
\newcommand{\vw}{\ensuremath{\vec{w}}}
\newcommand{\vx}{\ensuremath{\vec{x}}}
\newcommand{\vy}{\ensuremath{\vec{y}}}
\newcommand{\vb}{\ensuremath{\vec{b}}}
\newcommand{\vo}{\ensuremath{\vec{0}}}
\newcommand{\va}{\ensuremath{\vec{a}}}
\newcommand{\ve}{\ensuremath{\vec{e}}}

\newcommand{\R}{\mathbb{R}}
\newcommand{\Z}{\mathbb{Z}}
\newcommand{\C}{\mathbb{C}}
\newcommand{\N}{\mathbb{N}}
\newcommand{\Q}{\mathbb{Q}}

%%%%%%%%%%%%%%%%%%%%%%%%%%%%%%%%%%%%%%%%%%%%%%%%%%%%%%%%%%%

\begin{document}
	\begin{center}
		\Large{Math 251W: Foundations of Advanced Mathematics}

		\normalsize Portfolio problems from section 3.4

		\hfill Name: Carmen Beltran, Brendan Butler, Teddy Wachtler
		\hrule
		\vspace{0.4cm}
		\vspace{0.75cm}
	\end{center}

%%%%%%%%%%%%%%%%%%%%%%%%%%%%%%%%%%%%%%%%%%%%%%%%%%%%%%%%%%%

	\noindent Problem 3.3.8
	\vspace{0.3cm}

	\begin{enumerate}
		\item[v]
		\underline{proposition:}
		Given sets $A$, $B$, $C$, and $D$, $(A \cap B) \times (C \cap D) = (A \times C) \cap (B \times D)$.

		\vspace{.2cm}

		\fbox{proof} (Direct)
		Let $A$, $B$, $C$ and $D$ be arbitrary sets.

		Let $(x_1, x_2)$ be an arbitrary element of $(A \cap B) \times (C \cap D)$.
		By the definition of the cartesian product, we see $x_1 \in (A \cap B)$ and $x_2 \in (C \cap D)$.
		By the definition of set intersection, we see $x_1 \in A$, $x_1 \in B$, $x_2 \in C$, and $x_2 \in D$.
		By the properties of the cartesian product, since $x_1 \in A$ and $x_2 \in C$, we see that $(x_1, x_2) \in A \times C$, and since $x_1 \in B$ and $x_2 \in D$, we see that $(x_1, x_2) \in B \times D$.
		By the definition of set intersection, since $(x_1, x_2) \in A \times C$ and $(x_1, x_2) \in B \times D$, we see $(x_1, x_2) \in (A \times C) \cap (B \times D)$.
		Therefore, since $(x_1, x_2)$ was an arbitrary element of $(A \cap B) \times (C \cap D)$, we see $(A \cap B) \times (C \cap D) \subseteq (A \times C) \cap (B \times D)$.3

		Let $(y_1, y_2)$ be an arbitrary element of $(A \times C) \cap (B \times D)$.
		By the definition of set intersection, then [...].
		We see $(y_1, y_2) \in (A \cap B) \times (C \cap D)$.
		Therefore, since $(y_1, y_2)$ was an arbitrary element of $(A \times C) \cap (B \times D)$, we see $(A \times C) \cap (B \times D) \subseteq (A \cap B) \times (C \cap D)$.

		By the definition of set equality, we see that $(A \cap B) \times (C \cap D) = (A \times C) \cap (B \times D)$.
		Therefore, for all sets $A$, $B$, $C$, and $D$, $(A \cap B) \times (C \cap D) = (A \times C) \cap (B \times D)$.

		$\qed$.

		\vspace{.2cm}
	\end{enumerate}
	\vspace{1cm}

%%%%%%%%%%%%%%%%%%%%%%%%%%%%%%%%%%%%%%%%%%%%%%%%%%%%%%%%%%%

	\noindent Problem 3.4.3
	\vspace{0.3cm}

	\begin{enumerate}
		\item[vi]

		\underline{proposition:}
		If $I$ is a nonempty set, $\{A_i\}_{i \in I}$ is a family of sets indexed by $I$, and $B$ is a set, then $B - (\bigcap_{i\in I} A_i) = \bigcup_{i\in I}(B - A_i)$.

		\vspace{.3cm}

		\fbox{Proof} ( )

		\vspace{.2cm}
	\end{enumerate}
	\vspace{1cm}

%%%%%%%%%%%%%%%%%%%%%%%%%%%%%%%%%%%%%%%%%%%%%%%%%%%%%%%%%%%

	\noindent Problem 3.4.6
	\vspace{0.3cm}

	\begin{enumerate}
		\item[2]
		\underline{proposition:}
		If $I$ is a nonempty set, $\{A_i\}_{i \in I}$ is a family of sets indexed by $I$, and $B$ is a set, then $B\times(\bigcap_{i \in I} A_i) = \bigcap_{i \in I}(B \times A_i)$.

		\vspace{.3cm}

		\fbox{Proof} (Direct)



		\vspace{.2cm}
	\end{enumerate}
	\vspace{1cm}

%%%%%%%%%%%%%%%%%%%%%%%%%%%%%%%%%%%%%%%%%%%%%%%%%%%%%%%%%%%

\end{document}
